%
% Die Cauchy-Riemann-Differentialgleichungen
%
\subsection{Cauchy-Riemann-Differentialgleichung}
Die komplexe Ableitung $f'(z)$ einer komplexen Funktion $f(z)$
nach Definition~\ref{buch:holomorph:diffbar:def:komplexeableitung}
ist eine komplexe Zahl, mit der $\Delta z$ multipliziert werden muss.
Diese Multiplikation kann als Multiplikation der Matrix
\[
\begin{pmatrix}
\operatorname{Re}f'(z)
&
-\operatorname{Im}f'(z)
\\
\operatorname{Im}f'(z)
&
\phantom{-}\operatorname{Re}f'(z)
\end{pmatrix}
\qquad\text{mit}\qquad
\begin{pmatrix}
\operatorname{Re}\Delta z\\
\operatorname{Im}\Delta z
\end{pmatrix}
\]
betrachtet werden.

Andererseits zeigt \eqref{buch:holomorph:holomorph:2dersatz},
dass diese Matrix auch mit der Jacobi-Matrix der Funktion $f$
betrachtet als einer vektorwertigen Funktion von zwei reellen Variablen
übereinstimmen muss.
Es gilt also
\[
\begin{pmatrix}
\operatorname{Re}f'(z)
&
-\operatorname{Im}f'(z)
\\
\operatorname{Im}f'(z)
&
\phantom{-}\operatorname{Re}f'(z)
\end{pmatrix}
=
\Jf
=
\renewcommand{\arraystretch}{1.8}
\begin{pmatrix}
\displaystyle \frac{\partial u}{\partial x}(x,y)
&
\displaystyle \frac{\partial u}{\partial y}(x,y)
\\
\displaystyle \frac{\partial v}{\partial x}(x,y)
&
\displaystyle \frac{\partial v}{\partial y}(x,y)
\end{pmatrix}.
\]
Insbesondere gilt
\begin{align*}
\frac{\partial u}{\partial x}(x,y)
&=
\operatorname{Re}f'(z)
=
\frac{\partial v}{\partial y}(x,y)
&&\text{und}&
-\frac{\partial u}{\partial y}(x,y)
&=
\operatorname{Im}f'(z)
=
\frac{\partial v}{\partial x}(x,y)
\end{align*}
Dies sind zwei partielle Differentialgleichungen für Real- und Imaginärteil
einer holomorphen Funktionen.

\begin{satz}[Cauchy-Riemann-Differentialgleichungen]
\label{buch:holomorph:holomorph:satz:cr}
Realteil $u(x,y)=\operatorname{Re}f(x+iy)$ und
Imaginärteil $v(x,y)=\operatorname{Im}f(x+iy)$
einer in $U\subset \mathbb{R}$ holomorphen Funktion $f(z)$ erfüllen die
partiellen Differentialgleichungen
\begin{align}
\frac{\partial u}{\partial x}
&=
\frac{\partial v}{\partial y}
&&\text{und}&
\frac{\partial u}{\partial y}
&=
-
\frac{\partial v}{\partial x}.
\label{buch:holomorph:holomorph:eqn:cr}
\end{align}
Sie heissen die \emph{Cauchy-Riemann-Differentialgleichungen}.
\index{Cauchy-Riemann-Differentialgleichungen}%
\end{satz}

%
% Harmonische Funktionen
%
\subsubsection{Harmonische Funktionen}
Der Laplace-Operator einer Funktion $u(x_1,\dots,x_n)$ von $n$
\index{Laplace-Operator}%
reellen Variablen ist definiert als
\[
\Delta u
=
\frac{\partial^2u}{\partial x_1^2}
+
\dots
+
\frac{\partial^2u}{\partial x_n^2}.
\]
Der Laplace-Operator taucht in vielen Anwendungen in der Physik und
der Technik auf.
Dies hängt damit zusammen, dass er mehr oder weniger der einzige
partielle Differentialoperator zweiter Ordnung ist, der gegenüber
Drehungen des Raumes invariant ist.
Zum Beispiel erfüllt das Potential $u$ einer Ladungsverteilung $\varrho$ im
\index{Potential}%
\index{Ladungsverteilung}%
Vakuum die Laplace-Glei\-chung
\[
\Delta u = 4\pi \varrho.
\]
Entsprechend häufig ist man in den Anwendungen mit der Frage konfrontiert,
ob eine Funktion $u$ eine Lösung der homogenen partiellen Differentialgleichung
$\Delta u=0$ ist.

\begin{definition}[harmonische Funktion]
Eine \emph{harmonische Funktion} ist eine Funktion
$u\colon\mathbb{R}^n\to\mathbb{R}$, die $\Delta u = 0$ erfüllt.
\index{harmonische Funktion}%
\end{definition}

Harmonische Funktionen erfüllen also eine sehr spezielle partielle
Differentialgleichung.
Es ist auf den ersten Blick ganz und gar unklar, ob es überhaupt
solche Funktionen gibt.
Es zeigt sich aber, dass die Theorie der holomorphen Funktionen
dank des folgenden Satzes eine sehr grosse Zahl solcher Funktionen
liefert.

\begin{satz}
\label{buch:holomorph:holomorph:satz:harmonisch}
Real- und Imaginärteil einer holomorphen Funktion sind harmonische
Funktionen.
\end{satz}

\begin{proof}
Ableiten der ersten Cauchy-Riemann-Gleichung in
\eqref{buch:holomorph:holomorph:eqn:cr}
nach $x$ und der zweiten nach $y$ ergibt
\[
\frac{\partial^2 u}{\partial x^2}
=
\frac{\partial^2 u}{\partial y\,\partial x}
\quad\text{und}\quad
\frac{\partial^2 u}{\partial y^2}
=
-
\frac{\partial^2 v}{\partial x\,\partial y}
\quad\Rightarrow\quad
\frac{\partial^2 u}{\partial x^2}=-\frac{\partial^2 u}{\partial y^2}
\quad\Rightarrow\quad
\frac{\partial^2 u}{\partial x^2} + \frac{\partial^2 u}{\partial y^2}
=
\Delta u
=
0.
\]
Ebenso ergibt Ableiten der ersten Cauchy-Riemann-Differentialgleichung nach
$y$ und der zweiten nach $x$
\[
\frac{\partial^2 u}{\partial x\,\partial y}
=
\frac{\partial^2 v}{\partial y^2}
\quad\text{und}\quad
\frac{\partial^2 u}{\partial x\,\partial y}
=
-
\frac{\partial v}{\partial x^2}
\quad\Rightarrow\quad
\frac{\partial^2 v}{\partial y^2}
=
-
\frac{\partial^2 v}{\partial x^2}
\quad\Rightarrow\quad
\frac{\partial^2 v}{\partial x^2}
+
\frac{\partial^2 v}{\partial y^2}
=
\Delta v
=
0.
\]
Somit sind sowohl der Real- wie der Imaginärteil von $f$ harmonisch.
\end{proof}

\begin{beispiel}
\label{buch:holomorph:holomorph:bsp:zkholomorph}
Die Potenzfunktion $z\mapsto z^k$ ist holomorph, also müssen nach
Satz~\ref{buch:holomorph:holomorph:satz:harmonisch} Real- und
Imaginärteil harmonisch sein.
Wir prüfen dies am Beispiel $k=3$ nach, in diesem Fall ist
\begin{equation*}
(x+iy)^3
=
x^3 +3ix^2y - 3xy^2 -iy^3
\qquad\Rightarrow\qquad
\left\{
\begin{aligned}
u(x,y)
&=
x^3-3xy^2
\\
v(x,y)
&=
3x^2y-y^3.
\end{aligned}
\right.
\end{equation*}
Wir berechnen die zweiten partiellen Ableitungen für beide
Funktionen und erhalten
\begin{align*}
&\left.
\begin{aligned}
\frac{\partial u}{\partial x} &= 3x^2 -3y^2,
&
\frac{\partial^2 u}{\partial x^2} &= 6x,
\\
\frac{\partial u}{\partial y} &= -6xy,
&
\frac{\partial^2 u}{\partial y^2} &= -6x
\end{aligned}
\quad
\right\}
&&\Rightarrow&
\Delta u
&=
\frac{\partial^2 u}{\partial x^2}
+
\frac{\partial^2 u}{\partial y^2}
=
6x-6x = 0
\\
&
\left.\begin{aligned}
\frac{\partial v}{\partial x} &= 6xy,
&
\frac{\partial^2 v}{\partial x^2} &= 6y,
\\
\frac{\partial v}{\partial y} &= 3x^2-3y^2,
&
\frac{\partial v}{\partial y^2} &= -6y
\end{aligned}
\quad
\right\}
&&\Rightarrow&
\Delta v
&=
\frac{\partial^2 v}{\partial x^2}
+
\frac{\partial^2 v}{\partial y^2}
=
6y-6y
=
0.
\end{align*}
Beide sind also harmonisch.
\end{beispiel}

%
% Reelle Funktion zu einer holomorphen Funktion erweitern
%
\subsubsection{Reelle Funktion zu einer holomorphen Funktion erweitern}
Der Realteil einer holomorphen Funktion muss nach
Satz~\ref{buch:holomorph:holomorph:satz:harmonisch}
eine harmonische Funktion sein.
Es stellt sich damit auch die Frage, ob sich eine beliebige
reelle harmonische Funktion zu einer holomorphen Funktion 
erweitern lässt.

\begin{satz}
\label{buch:holomorph:holomorph:satz:utov}
Sei $u(x,y)$ eine harmonische reelle Funktion in einer offenen Umgebung
eines Punktes $z_0 = x_0+iy_0$.
Dann gibt es eine Funktion $v(x,y)$ in einer möglicherweise kleinern
Umgebung des Punktes derart, dass $f=u+iy$ eine holomorphe Funktion ist.
Die Funktion $v(x,y)$ ist bis auf eine reelle Konstante eindeutig
bestimmt.
\end{satz}

\input{chapters/020-holomorph/fig/fig-imaginaerteil.tex}%

\begin{proof}
Damit $f$ holomorph ist, muss $v$ so konstruiert werden, dass die
Cauchy-Riemann-Differentialgleichugen gelten.
Diese legen nur die Ableitungen fest, bleiben also erhalten, wenn man
zu $v$ eine beliebige Konstante hinzuaddiert.
Die Funktion $v$ ist also ist also nur bis auf eine Konstante
festgelegt.
Wir können daher $v(x_0,y_0)=v_0\in\mathbb{R}$ beliebig wählen.

Die Cauchy-Riemann-Differentialgleichungen legen die Ableitungen 
von $v$ nach $x$ und $y$ fest:
\begin{align*}
\frac{\partial v}{\partial x}(x,y)
&=
-\frac{\partial u}{\partial y}(x,y)
&&\text{und}&
\frac{\partial v}{\partial y}(x,y)
&=
 \frac{\partial u}{\partial x}(x,y)
\end{align*}
Durch Integration dieser beiden Gleichungen (siehe auch
Abbildung~\ref{buch:holomorph:holomorph:fig:imaginaerteil}) können die Werte 
\begin{align}
v(x,y_0)
&=
v_0
+
{\color{darkred}
\int_{x_0}^x
\frac{\partial v}{\partial x}(\xi,y_0)
\,d\xi}
=
v_0
{\color{darkred}\mathstrut-
\int_{x_0}^x
\frac{\partial u}{\partial y}(\xi,y_0)
\,d\xi}
\notag
\intertext{von $v$ berechnet werden.
Durch Integration der zweiten Gleichung kann jetzt auch der Wert}
v(x,y)
&=
v(x,y_0)
+
{\color{blue}
\int_{y_0}^y
\frac{\partial v}{\partial y}(x,\eta)\
\,d\eta
}
=
v(x,y_0) +
{\color{blue}
\int_{y_0}^y
\frac{\partial u}{\partial x}(x,\eta)
\,d\eta
}
\notag
\\
&=
v_0
{\color{darkred}\mathstrut
-
\int_{x_0}^x \frac{\partial u}{\partial y}(\xi,y_0)\,d\xi
}
+
{\color{blue}
\int_{y_0}^y \frac{\partial u}{\partial x}(x,\eta)\,d\eta
}
.
\label{buch:holomorph:holomorph:eqn:v}
\end{align}
bestimmt werden.

Die Formel~\eqref{buch:holomorph:holomorph:eqn:v} ist wohldefiniert,
solange die beiden Strecken von $x_0+iy_0$ nach $x+iy_0$ und
von $x+iy_0$ nach $x+iy$ vollständig in $U$ enthalten sind.
Nur dann ist der Integrand für beide Integrale definiert.

Die Formel~\eqref{buch:holomorph:holomorph:eqn:v} kann aber nur dann
der Imaginärteil einer holomorphen Funktion sein, wenn für diese
Funktion $v(x,y)$ tatsächlich die Cauchy-Riemann-Differentialgleichun\-gen
erfüllt.
Dazu müssen die Ableitungen nach $x$ und $y$ bestimmt werden:
\begin{align}
\frac{\partial v}{\partial y}(x,y)
&=
-\frac{\partial u}{\partial y}(x,y_0)
+
\int_{y_0}^y
\frac{\partial^2u}{\partial x^2}(x,\eta)
\,d\eta
\notag
\\
&=
-\frac{\partial u}{\partial y}(x,y_0)
-
\int_{y_0}^y 
\frac{\partial^2 u}{\partial y^2}(x,\eta)
\,d\eta
\label{buch:holomorph:holomorph:eqn:ulaplace}
\\
&=
-\frac{\partial u}{\partial y}(x,y_0)
-
\biggl[
\frac{\partial u}{\partial y}(x,\eta)
\biggr]_{y_0}^y 
\notag
\\
&=
-
\frac{\partial u}{\partial y}(x,y_0)
-
\frac{\partial u}{\partial y}(x,y)
+
\frac{\partial u}{\partial y}(x,y_0)
=
-\frac{\partial u}{\partial y}(x,y),
\notag
\\
\frac{\partial v}{\partial x}(x,y)
&=
\frac{\partial u}{\partial x}(x,y).
\notag
\end{align}
Im Schritt~\eqref{buch:holomorph:holomorph:eqn:ulaplace} haben wir 
$\Delta u=0$
in der Form
\[
\Delta u = 0
\quad\Rightarrow\quad
\frac{\partial^2u}{\partial x^2}
+
\frac{\partial^2u}{\partial y^2}
=
0
\quad\Rightarrow\quad
\frac{\partial^2u}{\partial x^2}
=
-
\frac{\partial^2u}{\partial y^2}
\]
verwendet.
Die Cauchy-Riemann-Differentialgleichungen sind also tatsächlich erfüllt
und damit ist $f(x+iy) = u(x,y)+ iv(x,y)$ eine holomorphe Funktion.

Der Wert $v_0=v(x_0,y_0)$ wird durch die
Cauchy-Riemann-Differentialgleichungen nicht festgelegt, er kann
beliebig gewählt werden.
Damit ist auch die Eindeutigkeit des Imaginärteils bis auf eine
Konstante gezeigt.
\end{proof}

\begin{beispiel}
Wir nehmen das Beispiel~\ref{buch:holomorph:holomorph:bsp:zkholomorph}
wieder auf und bestimmen den Imaginärteil einer holomorphen Funktion
mit dem Realteil $u(x,y)=x^3-3xy^2$.
Die Formel~\eqref{buch:holomorph:holomorph:eqn:v}
aus dem Beweis von Satz~\ref{buch:holomorph:holomorph:satz:utov}
sagt, dass dazu die partiellen Ableitungen von $u$ berechnet werden
müssen, sie sind
\begin{align*}
\frac{\partial u}{\partial x}
&=
3x^2-3y^2
&&\text{und}&
\frac{\partial u}{\partial y}
&=
-6xy.
\end{align*}
Anwendung der Formel~\eqref{buch:holomorph:holomorph:eqn:v}
liefert jetzt
\begin{align*}
v(x,y)
&=
v_0
-
\int_{x_0}^x -6\xi y_0\,d\xi
+
\int_{y_0}^y 3x^2-3\eta^2
\,d\eta
\\
&=
v_0
-
\biggl[
-3\xi^2y_0
\biggr]_{x_0}^x
+
\biggl[
3x^2\eta-\eta^3
\biggr]_{y_0}^y
\\
&=
v_0
+
3x^2y_0-3x_0^2y_0
+
3x^2y-y^3
-
3x^2y_0 + y_0^3
\\
&=
(v_0
+ y_0^3
-3x_0^2y_0
)
+
3x^2y
-
y^3.
\end{align*}
Bis auf den Klammerausdruck, der eine Konstante ist, wird also der
Ausdruck $3x^2y-y^3$ gefunden, der auch schon in
Beispiel~\ref{buch:holomorph:holomorph:bsp:zkholomorph}
aufgetreten ist.
\end{beispiel}

Der Satz funktioniert natürlich auch für einen vorgegebenen
Imaginärteil $v$.
Sucht man eine Funktion $f(z)$ mit diesem Imaginärteil, dann
hat $g(z)=-if(z)$ die Funktion $v$ als Realteil.
Nach Satz~\ref{buch:holomorph:holomorph:satz:utov} lässt sich
dann ein geeigneter Imaginärteil für $g(z)$ finden.
$f(z)=ig(z)$ ist schliesslich die gesuchte holomorphe Funktion.
